\documentclass[11pt,a4paper]{article}

% ---------------------------------------------------------
% PREAMBLE (stable, warning-free, Overleaf-safe)
% ---------------------------------------------------------
\usepackage[utf8]{inputenc}
\usepackage[T1]{fontenc}
\usepackage{lmodern}
\usepackage{amsmath, amssymb, amsfonts}
\usepackage{bm}
\usepackage{geometry}
\usepackage{graphicx}
\usepackage{hyperref}
\usepackage{cite}
\usepackage{authblk}
\usepackage{physics}
\usepackage{mathtools}

\geometry{margin=2.4cm}

\title{\textbf{Companion LXXXIV: Joint Constraints on $\alpha$ Using Multi-Scale Topological Transport and Minkowski Functionals in the Relaxation-Driven Cyclic Cosmology}}
\author{Michael Lehmann \\ \texttt{mi.lehmann@gmx.de}}
\date{April 2026}

\begin{document}
\maketitle

\begin{abstract}
Topological transport statistics and Minkowski functionals provide complementary 
probes of the large-scale structure encoded in weak-lensing convergence fields. 
Building on the transport-based parameter inference developed in previous 
Companions, this work constructs joint constraints on the relaxation parameter 
$\alpha$ in the Relaxation-Driven Cyclic Cosmology (RDCC) using a combined 
topology–morphology likelihood. Transport distances between persistence-diagram 
ensembles capture non-Gaussian topological deformations, while Minkowski 
functionals quantify morphological changes in the convergence field. We derive a 
joint Fisher matrix that incorporates cross-covariances between these statistics 
and identify the angular scales and topological features that contribute most 
strongly to constraining $\alpha$. This Companion provides the first unified 
multi-probe framework for RDCC parameter inference from Euclid-like surveys.
\end{abstract}
% ---------------------------------------------------------
\section{Introduction}

Weak-lensing convergence maps contain rich non-Gaussian information about the 
large-scale structure of the Universe. Persistence diagrams capture the topology 
of these maps, while Minkowski functionals quantify their morphology. Previous 
Companions developed a hierarchy of transport-based topological statistics and 
constructed Fisher forecasts for the RDCC relaxation parameter $\alpha$.

The present Companion unifies these approaches by combining topological transport 
statistics with Minkowski functionals in a joint likelihood. This multi-probe 
framework leverages the complementary sensitivity of topology and morphology to 
RDCC-induced deformations, yielding significantly improved constraints on 
$\alpha$.
% ---------------------------------------------------------
\section{Topological Transport Statistics}

Let $\mathcal{T}_{\ell}(\alpha)$ denote the scale-resolved transport discrepancy 
between $\Lambda$CDM and RDCC persistence diagrams, as defined in Companion 
LXXXII. These statistics capture:

\begin{itemize}
    \item shifts in birth–death structure,
    \item deformation of loop and void persistence,
    \item non-Gaussian topological signatures of RDCC.
\end{itemize}

Their sensitivity to $\alpha$ is encoded in the derivative 
$\partial \mathcal{T}_{\ell} / \partial \alpha$.
% ---------------------------------------------------------
\section{Minkowski Functionals as Morphological Probes}

For a convergence field $\kappa(\theta)$, the three Minkowski functionals in 2D 
are:

\begin{align}
    V_{0}(\nu) &= \text{area fraction above threshold}, \\
    V_{1}(\nu) &= \text{boundary length}, \\
    V_{2}(\nu) &= \text{Euler characteristic}.
\end{align}

RDCC modifies these functionals through the scale-dependent suppression of the 
gravitational potential. The derivatives 
$\partial V_{i}(\nu) / \partial \alpha$ quantify their sensitivity to $\alpha$.
% ---------------------------------------------------------
\section{Joint Likelihood and Fisher Matrix}

We define the combined data vector



\[
    \mathbf{d}(\alpha)
    =
    \left\{
        \mathcal{T}_{\ell}(\alpha),
        V_{0}(\nu),
        V_{1}(\nu),
        V_{2}(\nu)
    \right\}.
\]



The joint Fisher matrix is

\begin{equation}
    F_{\alpha\alpha}
    =
    \left(
        \frac{\partial \mathbf{d}}{\partial \alpha}
    \right)^{\!\!T}
    \mathbf{C}^{-1}
    \left(
        \frac{\partial \mathbf{d}}{\partial \alpha}
    \right),
\end{equation}

where $\mathbf{C}$ includes:

\begin{itemize}
    \item auto-covariances of transport statistics,
    \item auto-covariances of Minkowski functionals,
    \item cross-covariances between topology and morphology.
\end{itemize}

The forecasted constraint is $\sigma(\alpha) = F_{\alpha\alpha}^{-1/2}$.
% ---------------------------------------------------------
\section{Complementarity of Topology and Morphology}

Transport statistics and Minkowski functionals probe different aspects of the 
convergence field:

\begin{itemize}
    \item Transport statistics are sensitive to the \emph{connectivity structure} 
          of peaks, filaments, and voids.
    \item Minkowski functionals probe \emph{shape, area, and curvature}.
\end{itemize}

RDCC induces:

\begin{itemize}
    \item topological shifts in loop and void persistence,
    \item morphological changes in peak abundance and curvature,
    \item scale-dependent deformations that appear differently in both probes.
\end{itemize}

The joint analysis captures these complementary signatures.
% ---------------------------------------------------------
\section{Conclusion}

Topological transport statistics and Minkowski functionals provide complementary 
probes of RDCC-induced deformations in weak-lensing convergence fields. By 
combining them in a joint Fisher analysis, we obtain significantly improved 
constraints on the relaxation parameter $\alpha$. This multi-probe framework 
bridges the gap between geometric transport theory and observational cosmology, 
and is directly applicable to Euclid-like surveys.

\bibliographystyle{apsrev4-2}
\nocite{*}
\bibliography{references}


\end{document}
